\chapter{PROBLEM FORMULATION AND RESEARCH HYPOTHESIS}\label{ch:problem}

\begin{doublespacing}

This chapter formalizes the search problem, defines the metric used
throughout the thesis, introduces the sortie model for energy
constraints, and states the research hypotheses that the remainder of
the thesis tests.

\section{Search Instance}

\begin{definition}[Search Instance]
A search instance $\mathcal{I} = (\mathcal{N}, \mathbf{p}, \mathcal{B}, \tau)$
consists of (i) a set $\mathcal{N}$ of $n$ candidate sites with positions
$x_i \in [0, L]^2$, where $L$ is the side length of the square search
arena; (ii) a vector of prior probabilities
$\mathbf{p} = (p_1, \ldots, p_n)$ satisfying $p_i \geq 0$ and
$\sum_{i=1}^n p_i = 1$; (iii) a set $\mathcal{B}$ of $R$ robot base
stations at positions $x_0^r \in [0, L]^2$ for $r = 1, \ldots, R$; and
(iv) a target site $\tau \in \mathcal{N}$ drawn according to the prior
$\mathbf{p}$.
\end{definition}

Distances are Euclidean. The position of the target is written
$x_\tau$. A robot inspects a site by physically visiting it; a
visit reveals whether the target is present with certainty (the
perfect-sensor assumption, discussed in Chapter~VI as a limitation to
relax in future work). After each search round, robots return to base,
share their observations---specifically, the list of sites they
visited and did not find the target at---and update the joint prior
by Bayesian renormalization. For a set $\mathcal{V}$ of visited sites
from the current round, each visited site gets $p_i \leftarrow 0$ and
the remaining probabilities are renormalized:
\[
  p_i' = \frac{p_i}{1 - \sum_{j \in \mathcal{V}} p_j}, \quad i \notin \mathcal{V}.
\]
The updated prior is the input to the next round's auction.

\section{The Finder Competitive Ratio}

Competitive analysis requires a reference cost. We use the
\emph{optimal distance}, the distance the closest robot would need to
travel if it knew the target's location in advance:
\[
  D_{\mathrm{opt}} = \min_{r \in \{1,\ldots,R\}} d(x_0^r, x_\tau).
\]

\begin{definition}[Finder Competitive Ratio]
If robot $f$ is the robot that first inspects $\tau$, and $D_f$ is its
total travel distance up to the moment of discovery, then the
\emph{finder competitive ratio} (FCR) of the algorithm on instance
$\mathcal{I}$ is
\[
  \mathrm{FCR}(\mathcal{I}) = \frac{D_f}{D_{\mathrm{opt}}}.
\]
\end{definition}

FCR is the natural metric for this problem because it factors out the
underlying geometric difficulty of the instance. A team that happens
to be assigned a target near one of the bases should not be rewarded
for ``finding quickly''; FCR normalizes that away. An FCR of three
means the finding robot traveled three times the distance an omniscient
agent would have.

\section{Sortie Model}

In energy-constrained models, a robot begins a search at its base,
visits one or more sites, and returns to base to recharge. We call
each such excursion a \emph{sortie}.

\begin{definition}[Energy-Feasible Sortie]
A sortie for robot $r$ that visits sites $s_1, s_2, \ldots, s_k$ in
order is \emph{energy-feasible} under budget $E$ if
\[
  d(x_0^r, x_{s_1}) + \sum_{j=1}^{k-1} d(x_{s_j}, x_{s_{j+1}}) + d(x_{s_k}, x_0^r) \leq E.
\]
If the target is found mid-tour at site $s_j$, the finder stops
immediately and pays only $d(x_0^r, x_{s_1}) + \sum_{\ell=1}^{j-1}
d(x_{s_\ell}, x_{s_{\ell+1}}) + d(x_{s_{j-1}}, x_{s_j})$ (no return leg).
\end{definition}

We require every instance to be solvable: the target must be reachable
by at least one robot with a single-site sortie, that is,
$\min_r 2 \, d(x_0^r, x_\tau) \leq E$.

\section{Geometric Constants}

The bounds derived in Chapter~IV are expressed in terms of two
geometric quantities.

\begin{enumerate}
\item The expected distance between two independent uniform-random
points in the square $[0,L]^2$ is
$d_{\mathrm{avg}} \approx 0.5214\,L$~\cite{Robbins1978}. This is used
to estimate the expected cost of base-to-site legs.
\item The characteristic inter-site spacing
$d_{\mathrm{hop}} = L / \sqrt{n}$ arises from the regular-grid
approximation in which $n$ sites are placed in a square of side $L$;
each site then occupies a cell of area $L^2/n$, whose side length is
$L/\sqrt{n}$.
\end{enumerate}

\section{Research Hypotheses}

The research question stated in the introduction---how should robots
coordinate to minimize travel distance under energy and communication
constraints---decomposes into four testable hypotheses. Each is stated
below as a falsifiable claim, and each is evaluated in Chapter~V both
theoretically (bounds) and experimentally (Monte Carlo).

\begin{description}
\item[H1. Communication dominates fleet size.] An auction-coordinated
team of $R$ robots achieves substantially lower FCR than a team of $R$
robots that chooses sites independently. The hypothesis is that the
improvement factor is large (of order $5\times$ or greater) at the
default parameter setting.

\item[H2. Probability-awareness dominates centralization.] A
distributed auction with bids $p_i / d$ achieves lower FCR than a
centralized Hungarian assignment that minimizes total distance without
using probability information. Moreover, if the Hungarian is given
access to $p_i$-weighted costs, its performance matches the distributed
auction's, so that the gap is attributable to the bid function rather
than to the allocation mechanism.

\item[H3. Multi-site sorties recover most of the energy penalty.]
A multi-sortie planner that chains several sites per excursion closes
the majority of the FCR gap between an energy-constrained single-sortie
algorithm and an unconstrained-energy baseline.

\item[H4. A quadratic distance penalty is empirically optimal under
energy constraints.] Of the bid functions
$\{p, d^{-1}, p\,e^{-d/E}, p/d, p/d^2\}$, the function $p/d^2$ achieves
the lowest FCR on the default parameter setting. The hypothesized
mechanism is a cost-foreclosure argument: committing to a distant
anchor site wastes budget that would otherwise support additional
chain extensions.
\end{description}

\end{doublespacing}
